\chapter*{How to study this book}
\addcontentsline{toc}{chapter}{How to study this book}
This book assumes no previous optics, lithography, aberration theory, or
machine-learning knowledge. Basic algebra is useful. Calculus, complex
numbers, Fourier analysis, and matrix ideas are introduced as they become
necessary. The aim is to make every technical step in Chen and colleagues'
paper \cite{chen} understandable, including the steps compressed into
phrases such as rigorous simulation, Zernike coefficient, process variation
band, and inverse sensitivity analysis.

The research paper contains only three numbered equations. Those equations
are not the whole theory behind the work. The missing chain is
\[
\text{field}\ \longrightarrow\ \text{mask diffraction}\
\longrightarrow\ \text{pupil phase}\ \longrightarrow\ \text{image}\
\longrightarrow\ \text{edges}\ \longrightarrow\ \text{constraints}.
\]
The neural network is an approximate way to travel backward along part of
this chain. Learning the forward physics first makes the backward problem
much easier to understand.

The explanations are original teaching material. Source equations are
rederived in explicit conventions rather than copied without context.
Numerical demonstrations labelled \emph{teaching examples} are calculations
made for this book. They are not recovered data from the authors' commercial
simulations. In particular, no trained version of their neural network,
mask-simulator settings file, or proprietary dataset has been supplied.

\section*{A practical study sequence}
\begin{enumerate}
\item Read Chapters \ref{ch:story}--\ref{ch:waves} with a notebook. You
should be able to convert nanometres of optical path into radians of phase
and explain why field amplitudes interfere.
\item Study Chapters \ref{ch:diffraction}--\ref{ch:zernike}. Draw a pupil,
an image plane, a defocused wavefront, and a Zernike mode. Do the small
integrals yourself.
\item Work through Chapters \ref{ch:partial}--\ref{ch:sensitivity}.
These connect optics to the three measured quantities in the paper.
\item Read the data and learning chapters, then open the paper alongside
Chapter \ref{ch:walkthrough}. Every numbered equation, table, and figure
has a corresponding explanation.
\item Finish the budget and laboratory chapters and the solved exercises.
Try each exercise before looking at its answer.
\end{enumerate}
Do not try to memorize 25 polynomial labels. Understand the pupil function,
the unit convention, and the map from phase to edges. Those ideas remain
useful when the numbering convention or lithographic pattern changes.

\section*{What was actually available in the source files}
\begin{itemize}
\item The paper is complete: 10 PDF pages, journal pages 7270--7279.
Paper p.~3 means the third PDF page, journal page 7272.
\item The supplied \emph{Principles of Optics} scan is the fourth edition
of Born and Wolf \cite{bw}, with 859 PDF pages. Its printed pagination
and PDF pagination differ, and a central run is reversed. For example,
printed pages 464, 465, and 466 are PDF pages 555, 554, and 553.
Use the verified entry points in Appendix \ref{app:reading}, then follow
the printed page numbers.
\item The file named \texttt{Saleh \& Teich Foundations of Photonics.pdf}
actually contains \emph{Fundamentals of Photonics}, second edition
\cite{st}. It has 300 PDF pages and ends at printed page 278.
The useful chapters on ray optics, waves, Fourier optics,
electromagnetic optics, and polarization are present. Chapter 11,
\emph{Statistical Optics}, is listed in the contents but is not included.
This book therefore uses Born--Wolf for the missing partial-coherence
foundation and supplies its own complete derivations.
\end{itemize}

\section*{How to interpret statements in this book}
\textbf{Reported result} means something explicitly stated or displayed
in the paper. \textbf{Derived result} follows from equations and assumptions
written here. \textbf{Teaching example} uses a stated simplified model.
\textbf{Unresolved detail} means the publication does not supply enough
information to settle it. These distinctions matter because an attractive
graph is not a complete simulation specification.

Several publication issues deserve particular attention: the pattern tables
disagree, the output count and stated Zernike range disagree, a Table 4
column header is repeated, and the source-to-pupil conventions are incomplete.
The guide preserves the reported information and identifies the uncertainty.
It does not quietly choose a correction and present it as an author's fact.

This textbook concerns the supplied 2023 study. Its conclusions are read
in the context of its patterns, illumination, simulation conditions, and
validation data. It supplies the underlying education without treating a
simulation study as a complete production qualification.
